\chapter*{Resum}
\addcontentsline{toc}{chapter}{Resum}
Els videojocs d'estratègia complexos constitueixen un banc de proves exigent per a la recerca en Intel·ligència Artificial. En contrast amb jocs d'informació perfecta (com els escacs o el Go), els entorns amb informació parcial, accions simultànies i dinàmiques estocàstiques presenten reptes d'alta complexitat en predicció, modelització de l'oponent i planificació sota incertesa. Aquest Treball de Final de Màster desenvolupa, instrumenta i avalua de manera exhaustiva cinc paradigmes d'agents per a la presa de decisions en el format individual (\texttt{gen9randombattle}) del videojoc competitiu Pokémon, utilitzant el simulador de codi obert Pokémon Showdown i el marc d'integració \textit{poke-env}.

El problema es formalitza com un Procés de Decisió de Markov Parcialment Observable (POMDP) amb accions simultànies, caracteritzat per un espai d'estats i accions combinatòriament dens i múltiples fonts d'incertesa (composicions d'equip i moviments ocults, predicció de canvis i efectes aleatoris en el càlcul de dany). L'objectiu no és únicament aproximar una política que maximitzi la taxa de victòries sota pressupostos computacionals explícits, sinó també analitzar i comprendre els patrons de comportament dels agents en situacions crítiques: la gestió del risc, la predicció de canvis (\emph{switch prediction}), el control del ritme (\emph{tempo}) i l'ús estratègic de la Teracristal·lització, avaluats mitjançant telemetria detallada per torn i mètriques explicatives de decisió.

La proposta integra: (1) una seqüència estructurada de catorze agents heurístics (\texttt{v1}--\texttt{v14}) que incorporen incrementalment coneixement de domini (càlcul de dany, efectivitat de tipus, perills d'entrada, moviments de preparació, gestió d'estats alterats i Teracristal·lització); (2) variants de cerca adversària d'un torn (\emph{minimax} \texttt{v15}--\texttt{v17}) amb funcions d'avaluació heurístiques a les fulles; (3) cerca en arbre de Monte Carlo rasa (\emph{shallow} MCTS \texttt{v18}--\texttt{v20}) amb simulacions de cinc torns mitjançant \texttt{LocalSim}, comparant selecció UCB1 amb selecció PUCT guiada per priors heurístics; (4) models d'aprenentatge per imitació (\texttt{v21}--\texttt{v22}) basats en XGBoost entrenats sobre partides humanes d'alt nivell (Elo~$1800+$) per predir la decisió macro d'atacar o canviar; (5) un marc d'aprenentatge per reforç profund (\emph{MaskablePPO}) amb espai d'accions emmascarat i inicialització per clonatge conductual (\emph{Behavioral Cloning}); i (6) una infraestructura paral·lela i recuperable d'avaluació que instrumenta telemetria per torn i càlcul de rànquings Bradley-Terry.

L'avaluació empírica comprèn una matriu competitiva de 28 agents (més de 6,4 milions de batalles), proves específiques de DRL i la integració de models LLM (\textit{PokéChamp} i \textit{PokéLLMon}). Els resultats demostren que l'heurística experta (\texttt{v12}) lidera el rendiment superant el referent competitiu \textit{Abyssal}. Així mateix, l'anàlisi de telemetria revela que tant la cerca en arbre com els models apresos necessiten recolzar-se en el criteri heurístic de base per ser competitius en entorns amb informació parcial.

\vspace{0.5cm}
\noindent\textbf{Paraules clau:} Intel·ligència Artificial en jocs, presa de decisions sota incertesa, POMDP, Pokémon Showdown, agents heurístics, cerca adversària (Minimax), cerca en arbre de Monte Carlo (MCTS), aprenentatge per imitació (XGBoost), aprenentatge per reforç profund (MaskablePPO), models de llenguatge (LLM), telemetria, explicabilitat, model Bradley--Terry.
